%%%==========================================================================%%%
%%%              Packages and such like added by Norman Dunbar               %%%
%%%==========================================================================%%%

\usepackage{nameref}% To use \nameref{...}
\usepackage{wrapfig}% To wrap text around figures.
\usepackage{listings}% Guess!
\usepackage[toc,page]{appendix}% Gets a ``blank'' Appendix page.
\usepackage{longtable}% For tables extending over a page.
\usepackage[obeyspaces]{url}% Allow spaces in \path{}.
\usepackage{xurl}% Ensure \URLs break at /.:*-~'" characters.
\usepackage{xcolor}% Required for the setup of the listings.
\usepackage[LGR,T1]{fontenc}% To get Ohms symbols.
\usepackage{array}%To get table's centred and wrapped.
\usepackage{courier}% For listings, the default is dire!
\usepackage[unicode=true,
    bookmarks=true,
    bookmarksnumbered=false,
    bookmarksopen=false,
    breaklinks=true,
    pdfborder={0 0 0},
    pdfborderstyle={},
    backref=false,
    colorlinks=true]{hyperref}
\usepackage{multicol}
\usepackage{multirow}
\usepackage{pmboxdraw} % For `tree` diagrams.

\urlstyle{same}% Ensure URLs are in the same font as the text. Not in monotype!

\definecolor{ocre}{RGB}{243,102,25}
\definecolor{wwwDarkGreen}{HTML}{006400}
\definecolor{wwwDarkOrchid}{HTML}{9932CC}
\definecolor{wwwDarkOrange}{HTML}{FF8C00}


\input{arduinoLanguage.tex}% To get Arduino code highlighting

\lstdefinestyle{myArduino}{%% Define an Arduino style for use later %%
    language=Arduino,
    %% Add other words needing highlighting below %%
    morekeywords=[1]{},                  % [1] -> dark green
    morekeywords=[2]{FILE_WRITE},        % [2] -> light blue
    morekeywords=[3]{SD, File},          % [3] -> bold orange
    morekeywords=[4]{open, exists},      % [4] -> orange
    %% The lines below add a nifty box around the code %%
    frame=single,
    rulesepcolor=\color{arduinoBlue},
    numbers=left,
    basicstyle={\footnotesize\ttfamily}%To match lstlisting settings
}

\input{AVRAssemblyLanguage.tex}

\lstdefinestyle{AVRasm}{
    language=AVRAssembly,
    %% Add other words needing highlighting below %%
    morekeywords=[1]{},        % [1] -> dark blue - opcodes
    morekeywords=[2]{},        % [2] -> blue - registers & bit names
    morekeywords=[3]{},        % [3] -> bold orange - directives
    morekeywords=[4]{},        % [4] -> bold green - assembler functions
    %% The lines below add a nifty box around the code %%
    frame=single,
    rulesepcolor=\color{arduinoBlue},
    numbers=left,
    basicstyle={\footnotesize\ttfamily}%To match lstlisting settings
}

% Macros for Arduino Assembly Language.
\newcommand{\ardvers}{2.1.0}%VEesion of Arduino IDE in use.

\newcommand{\ohms}{\textgreek{W}}%Ohms symbol.
\newcommand{\app}[1]{\emph{#1}}%Application names.
\newcommand{\reg}[1]{\texttt{\uppercase{#1}}}%Register names.
\newcommand{\regmc}[1]{\texttt{#1}}%Register names.
\newcommand{\regbit}[1]{\texttt{\uppercase{#1}}}%Register bit names.
\newcommand{\bin}[1]{0b{#1}}%Binary numbers.
\newcommand{\hexa}[1]{0x\uppercase{{#1}}}%Hexadecimal numbers.
\newcommand{\varb}[1]{\texttt{#1}}%Variable names.
\newcommand{\code}[1]{\texttt{#1}}%Inline code segments.
\newcommand{\func}[1]{\texttt{#1}}%Procedure and function names.
\newcommand{\pin}[1]{\texttt{#1}}%Arduino/AVR pin names.
\newcommand{\unitname}[1]{\texttt{#1}}%ELF section names.
\newcommand{\externC}{\code{extern \textquotedbl{}C\textquotedbl}}
\newcommand{\shl}{\texttt{<}\texttt{<}~}% << in inline code.
\newcommand{\shr}{\texttt{>}\texttt{>}~}% >> in inline code.
\newcommand{\atmega}{ATmega328P}% Get capitalisation correct!
\newcommand{\lyxarrow}{$\rightarrow$}
\newcommand{\twi}{I$^2$C}


%%% Setup Hyperref configuration %%%
\hypersetup{pdftitle={Arduino Assembly Language},
    pdfauthor={Norman Dunbar},
    pdfsubject={Using AVR Assembly Language with the Arduino IDE.},
    pdfkeywords={Arduino, AVR, ATmega328P,  microcontroller, Uno, Duemilanove, Nano, Assembly Language}}


%%% Lyx specific additions %%%
\makeatletter

\DeclareRobustCommand{\greektext}{%
    \fontencoding{LGR}\selectfont\def\encodingdefault{LGR}}
\DeclareRobustCommand{\textgreek}[1]{\leavevmode{\greektext #1}}
\ProvideTextCommand{\~}{LGR}[1]{\char126#1}

\makeatother

%%% Settings for listings %%%
\lstset{breaklines=true,
    tabsize=4,
    basicstyle={\footnotesize\ttfamily}, %%% Overridden by SVMono class. :(
    basewidth={0.55em}, %%% Needed to fit listings on page %%%
    showspaces=false,
    showstringspaces=false,
    keepspaces=true,
    frame=single}

%%%==========================================================================%%%
%%%                End of Norman Dunbar's additions                          %%%
%%%==========================================================================%%%